\project{17}{The same driver twice: vendor layer against registers}{NUCLEO-H7A3ZI-Q with one shield}{Abstraction layers, code size and cycles, and the Rust variant}

\begin{keyfacts}
\item[Board] NUCLEO-H7A3ZI-Q with one shield on the Arduino header. Either sensor shield serves, because what is measured is the transaction and not the sensor
\item[Peripherals] I2C1 for the sensor read, GPIO for the toggle measurement, the DWT cycle counter, TIM2 and SysTick as the substituted timebases, USART3 for the console
\item[Toolchain] arm-none-eabi-gcc with CMake for the two C implementations, and a Rust cross toolchain for the third and fourth
\item[Operating system] Bare metal, plus one real-time framework that schedules on the interrupt controller
\item[Difficulty] 4 of 5
\item[Effort] 5 evenings of about four hours
\item[Deliverable] One sensor read implemented four ways that return byte-identical results, with code size and cycle counts measured for each and the method published alongside the numbers
\end{keyfacts}

\subsection*{Why this project}

Ask whether a vendor's hardware abstraction layer is worth its cost and you will
get a confident answer within seconds and a measurement within never. This is
the weakest-evidenced topic in this book. There is a great deal of opinion, one
piece of reproducible bench work, and no well-argued case from the
vendor-layer side at all, which is itself remarkable: the layer is used by an
enormous number of projects and nobody has published a defence of it with
numbers. The chapter says that plainly and then supplies the measurement.

The one measured comparison worth the name is an engineer's bench work on a
Cortex-M7 part from this family running at 550 MHz, which found the vendor's pin
toggle call running at a fraction of the speed of a direct register write. That
is a real result and it is cited here with the caveat this whole volume exists
for: it was taken on a different member of the family, with a different
reference manual, a different clock tree and a different bus structure. The
shape of the finding transfers; the numbers do not.

The reason to do this on a sensor read rather than on a pin toggle is that the
answer changes, and the change is the finding. A pin toggle is two instructions
of real work, so any overhead at all dominates it. A six-byte register read over
a bus at 400 kHz is about nine bytes of traffic at nine bit times each, which is
81 bit times, which is about 202 microseconds, which at 280 MHz is roughly
56,700 core cycles of waiting. A call path costing a few hundred cycles is a
fraction of one percent of that. Both statements are true at once, and a reader
who has only met the first one will make bad decisions with it. The useful
quantity is not the overhead but the ratio of overhead to work, and this chapter
measures that ratio on two operations that sit at opposite ends of it.

\begin{note}{What this chapter is allowed to claim}
That the two implementations return identical bytes, and what each one cost in
flash and in cycles on this part at a stated clock, a stated optimisation level
and a stated cache state. Not that one layer is better. Not that the result
transfers to another part. Not that the ratio found here holds for a different
operation. The value of a measurement is in its boundaries, and a chapter on
this subject that overreached would be adding to the problem it is about.
\end{note}

\subsection*{Prior art and what to reuse}

\begin{table}[H]\centering\small
\begin{tabular}{p{34mm}p{46mm}p{38mm}p{20mm}}
\toprule
Source & What it gives & What it does not & Licence \\
\midrule
One engineer's bench measurement of pin toggling on this silicon family & The only reproducible comparison of the vendor call against a direct register write, with a method described well enough to repeat & It was taken on a different member of this family at a different clock, so its numbers are the shape of the answer and not the answer & Reading only \\
The vendor's firmware package for this board & Parallel example trees at both abstraction levels for this exact board, which is the raw material this chapter measures & Its examples are under a licence tied to this vendor's silicon, so they are read and measured rather than copied into a portfolio & Mixed, flagged \\
The Rust hardware crate for this family & A feature that names this part, routing it through its twin's peripheral module with the correct reference manual selected. That is the right engineering call and it is explained here rather than hidden & No example for this exact board, and a different set of assumptions about ownership and initialisation & 0BSD \\
The Rust peripheral crate collection & A generated peripheral module for this part, with typed register access that is itself an abstraction worth measuring & Nothing above the register level & Apache-2.0 or MIT \\
The async framework & It names this exact package and pin count, which is a stronger statement of support than most projects make & No worked example for this part at the time of writing, only for its twin. That gap is worth filling & Apache-2.0 or MIT \\
The real-time framework & Scheduling built on the interrupt controller, which is an architecture-level mechanism, so it works here without anything part-specific & No peripheral support of its own, by design & Apache-2.0 or MIT \\
The flashing and debugging tool & A target manifest carrying this exact die & Nothing about abstraction layers & Apache-2.0 or MIT \\
\bottomrule
\end{tabular}
\caption{Prior art for chapter 17. The gap in this table is the interesting part: there is no entry for a well-argued case in favour of the vendor layer, because there is not one to cite.}
\end{table}

What is left to write is the measurement and its method, plus the two
implementations that have to agree before the measurement means anything. The
method is the deliverable. Anyone can produce a number; producing a number with
the clock, the optimisation level, the cache state, the repetition count and the
reported statistic all stated is what makes it usable by somebody else, and it
is exactly what the existing writing on this subject does not do.

\subsection*{Parts from the inventory}

\begin{table}[H]\centering\small
\begin{tabular}{p{40mm}p{66mm}p{40mm}}
\toprule
Part & Role & Interface \\
\midrule
NUCLEO-H7A3ZI-Q & The whole chapter. Its probe is also the flashing path for both toolchains & Micro USB to the host \\
One sensor shield & The target of the read. Either the four-sensor or the five-sensor shield serves, because the operation under test is a bus transaction & I2C1 on the Arduino header \\
A USB data cable & Power, programming and console & Micro USB \\
nRF-PPK2 & The marker pin recorder, and a charge figure per operation if you want one & IDD jumper and one digital input \\
Host PC & Two toolchains, two build systems, one comparison script & Windows with the Linux subsystem \\
\bottomrule
\end{tabular}
\caption{Inventory items used in chapter 17. Nothing is bought and nothing is soldered. One shield is on the board at a time.}
\end{table}

\subsection*{System architecture}

\diagram{c17_arch}{Four implementations of one operation, one equality harness and one set of instruments. The harness is what makes the four comparable; the instruments are what make the comparison a measurement rather than an impression.}

The figure has one property worth pointing out. The two C implementations and
the two Rust implementations form the same ladder with different names. In C the
rungs are the vendor layer, the low-level layer and raw register access. In Rust
they are the hardware abstraction crate, the generated peripheral crate and a
raw volatile write. The peripheral crate is not the bottom of the ladder even
though it looks like registers, because it is a typed, ownership-checked view of
them and that view has a cost that can be measured like any other.

\subsection*{Peripheral configuration}

\begin{table}[H]\centering\small
\begin{tabular}{p{22mm}p{26mm}p{24mm}p{30mm}p{30mm}}
\toprule
Peripheral & Mode & Clock source & Pins and function & Interrupt and transfers \\
\midrule
I2C1 & Fast mode, 400 kHz & Peripheral bus, source selected in RCC & D14 and D15 on the Arduino header. Confirm in MB1363 & Blocking in all four implementations, so that only the layer differs \\
GPIO output & Push-pull, maximum speed & Peripheral bus & One free Zio pin, toggled for the second operation & None \\
GPIO output & Push-pull marker & Peripheral bus & One free Zio pin into a PPK2 digital input & None \\
DWT cycle counter & Free running & Core clock & None & The primary instrument \\
TIM2 & Free running, 32-bit, 1 MHz & Peripheral bus & None & The first substituted timebase \\
SysTick & Free running, reload at maximum & Core clock or its divided form & None & The second substituted timebase \\
USART3 & Asynchronous, 115200 8N1 & Peripheral bus & The console pins of chapter 1 & Polled \\
\bottomrule
\end{tabular}
\caption{Peripheral configuration. Every implementation uses the blocking form of the transfer, because comparing a blocking vendor call against an interrupt-driven register implementation would measure the execution model rather than the abstraction.}
\end{table}

\subsection*{Wiring}

\diagram{c17_wiring}{The shield supplies a device to talk to and the toggled pin supplies the other end of the ratio. One marker wire reaches the instrument. Nothing else is placed by hand.}

3.3 V logic only, as everywhere in this book, and no 5 V signal into a pin. The
toggled pin drives nothing but the instrument's digital input, so there is no
load to argue about. There is no oscilloscope on this bench, which matters more
here than in most chapters: a toggle at these rates cannot be observed
externally with what is available, so the cycle counter is the instrument and
the instrument's own cost is subtracted from every reading. The substitute for a
scope is stated rather than implied, and the subtraction is in the harness.

\subsection*{Memory and timing budget}

\diagram{c17_mem}{The two registers a pin toggle can go through. One store to the set and reset register changes a pin with no read; a read, modify and write of the output data register is three accesses and has a window in it. The difference is atomicity first and speed second, and the layer you call decides which one you get.}

\begin{table}[H]\centering\small
\begin{tabular}{p{44mm}p{28mm}p{28mm}p{30mm}}
\toprule
Quantity & Budget & Measured & Margin \\
\midrule
Flash, vendor-layer read & 4 kB & not measured & not measured \\
Flash, register-level read & 1 kB & not measured & not measured \\
Cycles, vendor-layer read & 57,000 & not measured & not measured \\
Cycles, register-level read & 56,800 & not measured & not measured \\
Cycles, vendor-layer toggle & 40 & not measured & not measured \\
Cycles, register-level toggle & 4 & not measured & not measured \\
Bus time for the read at 400 kHz & 202 \textmu{}s computed & not measured & not measured \\
Cost of reading the cycle counter & subtracted & not measured & not measured \\
\bottomrule
\end{tabular}
\caption{The budget table. The two read rows are deliberately close together and the two toggle rows are deliberately far apart, because that is the prediction this chapter exists to test. The bus figure is arithmetic from nine bytes at nine bit times and is marked computed rather than measured.}
\end{table}

Three things have to be stated with every cycle count on this part or the number
is not reusable. The core clock, because a cycle is not a time. The optimisation
level, because a call that is inlined at one level and not at another is a
different program. And the cache state, because this core has instruction and
data caches and a first run through cold code can cost many times what the
tenth run costs. The harness reports the first run and the median of the
remainder separately, and Chapter 19 explains why the two differ.

\subsection*{Firmware design (UML)}

\diagram{c17_uml}{One read, two call paths. The right-hand path is what the left-hand path does on your behalf, and reading the two side by side is the fastest way to see what the abstraction is actually for.}

What the vendor layer does that the register implementation does not is the
substance of the chapter, and it is not nothing. It checks that the peripheral
is in a state where the call is legal. It takes a lock, so that two callers
cannot interleave. It computes a deadline from a millisecond tick and abandons
the transfer if the deadline passes, which turns a stuck bus into an error code
rather than a stopped program. It inspects the error flags afterwards and maps
them onto a small set of return values. And it presents the same signature
whether the transfer is blocking, interrupt-driven or performed by the transfer
engine, which is worth a great deal when a design changes its mind.

The register implementation does none of that by default, which is why the
honest form of this chapter's question is not whether the abstraction is worth
its cost but which of those five services you actually need. A register-level
implementation that adds a timeout and error mapping, because it needs them,
converges on the vendor layer and its measured cost converges too. That
convergence is a result, and it is the part of this subject that opinion pieces
consistently miss.

\subsection*{Data flow (ASCII)}

\begin{asciiart}
  one operation, defined once
     |
     +--> impl 1: C, vendor layer ----------+
     +--> impl 2: C, registers -------------+
     +--> impl 3: Rust, peripheral crate ---+
     +--> impl 4: Rust, hardware crate -----+
                                            |
                                            v
  +-----------------------------------------------------------------------------------+
  | equality harness on the host: the four byte strings must be identical             |
  |   before any number anywhere in this chapter is recorded                          |
  +-----------------------------------------------------------------------------------+
                                            |
                                            v
  +-----------------------------------------------------------------------------------+
  | instruments: DWT cycles, first run and median reported                            |
  |   separately; arm-none-eabi-size per object; TIM2 and                             |
  |   SysTick as the substituted timebases; marker pin to PPK2                        |
  +-----------------------------------------------------------------------------------+
                                            |
                                            v
  +-----------------------------------------------------------------------------------+
  | one table: core clock, optimisation level, cache state, repetition                |
  |   count, median and maximum, in every row                                         |
  +-----------------------------------------------------------------------------------+
\end{asciiart}

\subsection*{Repository layout}

\begin{plaincode}
nucleo-h7a3-two-layers/
  c/
    CMakeLists.txt
    src/op_vendor.c          # the sensor read through the vendor layer
    src/op_registers.c       # the same read, written against the registers
    src/op_toggle.c          # both toggles, the other end of the ratio
    src/measure.c            # DWT, and the subtraction of its own cost
    src/timebase_dwt.c       # peripheral substitution, implementation 1
    src/timebase_tim2.c      # peripheral substitution, implementation 2
    src/timebase_systick.c   # peripheral substitution, implementation 3
  rust/
    Cargo.toml               # the crate features are the load-bearing part
    src/bin/pac.rs           # the read through the generated peripheral crate
    src/bin/hal.rs           # the read through the hardware crate
    src/bin/embassy.rs       # the worked example this part did not have
    memory.x                 # checked against RM0455, not inherited
  tools/
    equality.py              # four byte strings, one assertion
    collect.py  report.py    # the table, with its conditions in it
  docs/METHOD.md             # clock, optimisation, cache state, N, statistic
  README.md
\end{plaincode}

\subsection*{Steps}

\begin{steps}
\item \textbf{Define the operation before writing any implementation.} Write it
down in \texttt{docs/METHOD.md}: read six bytes from a named register of the
sensor at a named address, over I2C1 at 400 kHz, blocking, with the result
printed as twelve hexadecimal characters. Ambiguity here is what makes most
published comparisons unusable, because two implementations that do slightly
different work are not two implementations of one operation.

\item \textbf{Write the equality harness first.} It is ten lines and it is what
stops the chapter from comparing two different programs.
\begin{pycode}
results = {name: read_line(port, name) for name in
           ("vendor", "registers", "pac", "hal")}
assert len(set(results.values())) == 1, results
print("identical:", next(iter(results.values())))
\end{pycode}

\item \textbf{Implementation one, the vendor layer.} Use the blocking memory
read from the vendor's layer, configured by its own initialisation structure.
Keep it idiomatic: the point is to measure the layer as it is actually used, not
a stripped version of it that nobody writes.
\begin{ccode}
/* The vendor's own call. Note what it takes that the register version
   does not: a handle carrying state and a lock, and a timeout. */
if (HAL_I2C_Mem_Read(&hi2c1, DEV_ADDR << 1, REG_ADDR,
                     I2C_MEMADD_SIZE_8BIT, buf, 6, TIMEOUT_MS) != HAL_OK) {
    return -1;
}
\end{ccode}

\item \textbf{Implementation two, the registers.} Write the transfer against the
peripheral's own registers. One register on this peripheral carries the target
address, the direction, the byte count, the start condition and the automatic
end condition together, which is the single clearest example in the part of what
an abstraction is hiding. Read the field positions in RM0455 before writing
them; this book does not print a register map it has not checked, and the
reference manual for this part is not the one most search results return.
\begin{ccode}
/* Phase one: write the register address. Phase two: repeated start and
   read six bytes. The transfer-configuration register carries the target
   address, the direction, the count, the start and the automatic end. */
i2c_write_cr2(DEV_ADDR, 1, I2C_WRITE, I2C_SOFTEND);
i2c_put_byte(REG_ADDR);
i2c_wait_tc();                       /* transfer complete, not stop */
i2c_write_cr2(DEV_ADDR, 6, I2C_READ, I2C_AUTOEND);
for (int i = 0; i < 6; i++) buf[i] = i2c_get_byte();
i2c_wait_stop_and_clear();
\end{ccode}
Everything the vendor layer does and this does not is now visible as an absence:
no state check, no lock, no deadline, no error mapping. Add them back one at a
time later and watch the cost converge.

\item \textbf{Measure code size, per object rather than per image.} An image
figure is dominated by whatever else is linked in. The honest comparison is the
size of the objects that implement the operation, plus whatever they drag in.
\begin{shellcode}
arm-none-eabi-size -t build/CMakeFiles/*.dir/src/op_vendor.c.obj
arm-none-eabi-size -t build/CMakeFiles/*.dir/src/op_registers.c.obj
arm-none-eabi-nm --size-sort -S build/two_layers.elf | tail -20
\end{shellcode}

\item \textbf{Measure cycles, and subtract the instrument.} Reading the cycle
counter costs cycles. Measure an empty interval first and subtract it from every
reading, then report the first run and the median of the rest separately.
\begin{ccode}
static uint32_t cyc_overhead;

void measure_calibrate(void)
{
    uint32_t best = UINT32_MAX;
    for (int i = 0; i < 64; i++) {
        const uint32_t a = DWT->CYCCNT;
        const uint32_t b = DWT->CYCCNT;
        if (b - a < best) best = b - a;
    }
    cyc_overhead = best;              /* subtracted from every reading */
}

uint32_t measure_run(void (*op)(void))
{
    const uint32_t a = DWT->CYCCNT;
    op();
    const uint32_t b = DWT->CYCCNT;
    return (b - a) - cyc_overhead;
}
\end{ccode}
Whether the cycle counter works on this part with no debugger attached is on
this volume's confirm list. Check it before relying on it, because if it needs a
probe the measurement is only available on a bench and not in the field.

\item \textbf{Do the same on a pin toggle and compute the ratio.} This is the
step that produces the chapter's finding. The toggle is where the published
comparison found a large difference, and the read is where the bus dominates.
Report both, and report overhead as a fraction of the operation rather than as a
number of cycles.
\begin{ccode}
static inline void toggle_registers(void)   /* one store, no read */
{
    GPIOx->BSRR = PIN_MASK;                  /* set   */
    GPIOx->BSRR = PIN_MASK << 16;            /* reset */
}
\end{ccode}
The register figure above is what that line is doing. A read, modify and write
of the output register would produce the same waveform and would not be atomic,
which is a correctness difference and not only a speed one.

\item \textbf{Substitute the peripheral, and keep the job the same.} This
chapter owns that axis for the whole book, so it is built here in full. The job
is to produce an elapsed interval; three peripherals can do it.
\begin{ccode}
/* One interface, three peripherals behind it. Selected at build time so
   that nothing is paid for at run time. */
typedef struct {
    void     (*start)(void);
    uint32_t (*ticks)(void);     /* free-running count */
    uint32_t hz;                 /* what a tick is worth */
    uint32_t width_bits;         /* where it wraps */
} timebase_t;

extern const timebase_t TB_DWT;      /* core clock, 32 bits, no prescale  */
extern const timebase_t TB_TIM2;     /* 1 MHz, 32 bits, survives more     */
extern const timebase_t TB_SYSTICK;  /* core clock, 24 bits, counts down  */
\end{ccode}
Then compare them on resolution, wrap period, the cost of a read, and
availability. At 280 MHz the cycle counter wraps in about 15.3 seconds, a 32-bit
timer at 1 MHz wraps in about 71.6 minutes, and a 24-bit down counter at the
core clock wraps in about 60 milliseconds, which makes the third one unusable
for anything longer than an interrupt handler and perfectly good inside one.

\item \textbf{The Rust variant, in full.} The crate features are the part of
this that matters and they are verified: the hardware crate has a feature naming
this part which routes to its twin's peripheral module and selects the right
reference manual, and the async framework names this exact package and pin
count.
\begin{plaincode}
[dependencies]
cortex-m      = "0.7"
cortex-m-rt   = "0.7"
# The feature names are the verified part of this file. Take the versions
# from your own lockfile rather than from a page in a book.
stm32h7xx-hal = { version = "0.16", features = ["stm32h7a3", "rt"] }
\end{plaincode}
That routing is worth a paragraph in your README rather than a shrug. The crate
is not pretending this part is its twin; it is stating that the peripheral
module and the reference manual selection of the twin are the correct ones for
this die, which is precisely the distinction this whole volume is about, made
correctly, in a build file.
\begin{rustcode}
// The same six bytes, through the generated peripheral crate and through
// the hardware crate. Two rungs of the same ladder, in one language.
let dp = pac::Peripherals::take().unwrap();
let mut buf = [0u8; 6];

// hardware crate: ownership, typed pins, a configured bus
let mut i2c = dp.I2C1.i2c((scl, sda), 400.kHz(), ccdr.peripheral.I2C1,
                          &ccdr.clocks);
i2c.write_read(DEV_ADDR, &[REG_ADDR], &mut buf).unwrap();
\end{rustcode}
Build and flash with the same probe the C build uses, which is the practical
reason this variant is cheap to add on this bench.
\begin{shellcode}
rustup target add thumbv7em-none-eabihf
cargo build --release --bin hal
probe-rs run --chip STM32H7A3ZITxQ target/thumbv7em-none-eabihf/release/hal
\end{shellcode}

\item \textbf{Write the worked example this part did not have.} The async
framework names this exact package and pin count and had, at the time of
writing, a worked example only for its twin. Writing one for this part is a
small piece of genuinely new work, it is exactly the kind of contribution that
makes a portfolio repository useful to somebody else, and it is worth offering
upstream. The real-time framework needs no such work, because it schedules on
the interrupt controller, which is an architecture-level mechanism rather than a
part-specific one, and that is worth one sentence in the README so a reader
understands why one of the two needed porting and the other did not.
\end{steps}

\subsection*{Build, flash and debug}

\diagram{c17_timing}{What the two toggles look like against the core clock, and where the cycles go in the two reads. The toggle is where the layer costs a multiple; the read is where it costs a rounding error. Both are true and the chapter reports both.}

\begin{shellcode}
cmake -B c/build -DCMAKE_TOOLCHAIN_FILE=cmake/arm-none-eabi.cmake -DOPT=O2
cmake --build c/build -j
probe-rs run --chip STM32H7A3ZITxQ c/build/two_layers.elf
python tools/equality.py --port /dev/ttyACM0
python tools/report.py --out docs/RESULTS.md
\end{shellcode}

Recovering a board that will not answer is Chapter 1's material. The failure
specific to this chapter is a comparison that looks fine and is not.

\begin{note}{When the two implementations return different bytes}
In order of likelihood: the address was shifted in one implementation and not in
the other, because the vendor layer expects the seven-bit address in the upper
bits of an eight-bit field and a register implementation usually does not; the
register implementation stopped the transfer after phase one instead of
producing a repeated start; the automatic end condition was configured in the
phase where it should not be, so the sensor saw two transactions rather than
one; the byte count was written to the wrong field; or the sensor auto-increments
its register pointer in one case and not in the other because a configuration
bit differs. Compare the first byte only, then two bytes, then six. The point at
which they diverge names the fault.
\end{note}

\subsection*{Verification and acceptance criteria}

\begin{itemize}
\item All four implementations return byte-identical results, proven by the
harness rather than by inspection, and the harness runs in the build.
\item The cost of the abstraction is a number with its conditions attached: core
clock, optimisation level, cache state, repetition count, and median together
with maximum. A number without those is not published.
\item The instrument's own cost is measured and subtracted, and the subtraction
is visible in the source rather than folded into a constant.
\item The first run and the median of the remainder are reported separately, and
the difference between them is acknowledged as a cache effect rather than
averaged away.
\item Code size is reported per object with what it drags in, not as a
difference between two image sizes.
\item The ratio of overhead to work is reported for both operations, and the
chapter's conclusion is stated as a function of that ratio rather than as a
verdict on the layer.
\item The timebase substitution is demonstrated on all three peripherals, with
the wrap period of each stated, and the cycle counter's dependence on a probe
resolved rather than assumed.
\item The Rust build flashes with the same probe as the C build and the
peripheral module routing is explained in the README in one paragraph.
\end{itemize}

\subsection*{Variants}

\begin{table}[H]\centering\small
\begin{tabular}{p{24mm}p{26mm}p{44mm}p{22mm}p{20mm}}
\toprule
Axis & Variant & What changes & Cost & Built in full in \\
\midrule
Language & Rust & The same read at two rungs of the ladder, with the hardware crate routing this part through its twin's peripheral module and selecting the right reference manual & A second toolchain, and a different model of ownership to learn & Here \\
Peripheral substitution & One job on three peripherals & The elapsed-interval job moved between the cycle counter, a 32-bit timer and the system tick, with resolution and wrap period compared & Each has a different wrap period and a different dependence on the debugger & Here \\
Peripheral substitution & The same read on another bus instance & The identical transfer through a different peripheral instance, which on this family can mean a different clock domain and therefore different behaviour in low-power modes & Pin availability while a shield is fitted, which is on the confirm list & Here, gated on the header map \\
Language & C++ & The same read behind a template that resolves the layer at compile time, with no run-time cost and a different set of arguments about readability & Toolchain settings to justify & Chapter 9 \\
Execution model & Interrupt and DMA forms & The vendor layer presents the same signature for all three, which is one of the services being paid for & A different measurement entirely & Chapters 3 and 4 \\
Operating system & An RTOS task & The lock the vendor layer takes starts to matter once two tasks can call the same peripheral & Priority inheritance and a mutex to size & Chapter 20 \\
\bottomrule
\end{tabular}
\caption{Variants for chapter 17. Two axes are built in full here, which is why this chapter is longer on implementation and shorter on narrative than its neighbours.}
\end{table}

\subsection*{Pitfalls}

\begin{itemize}
\item Comparing a blocking call against an interrupt-driven implementation. That
measures the execution model and attributes the result to the abstraction.
\item Quoting a cycle count without the clock, the optimisation level and the
cache state. The number is then unusable by anybody, including you, in six
months.
\item Measuring only the pin toggle. It is the operation where overhead
dominates, and generalising from it is how this subject acquired its confident
opinions.
\item Forgetting the address shift. The vendor layer and a register
implementation disagree about where the seven-bit address sits, and the symptom
is a device that does not acknowledge.
\item Reporting a mean over runs. The first run through cold code is a different
measurement from the tenth, and averaging them produces a number that describes
neither.
\item Taking a linker script or a memory layout from the better-known member of
this family for the Rust build. The same rule applies in every language:
\texttt{memory.x} is checked against RM0455.
\item Treating the generated peripheral crate as the bottom of the ladder. It is
a typed view of registers with its own cost, and measuring it as though it were
raw access hides a rung.
\item Concluding that the vendor layer is expensive. What was measured is that
it is expensive relative to an operation with almost no work in it, which is a
much narrower statement and the only one the data supports.
\end{itemize}

\subsection*{Best practices applied}

\begin{itemize}
\item The equality harness exists before either implementation, so the
comparison cannot quietly become a comparison of two different operations.
\item The instrument's own cost is measured and subtracted rather than assumed
negligible.
\item Every published number carries its conditions in the same table, which is
the single thing the existing writing on this subject does not do.
\item The absence of prior art is stated as a finding rather than filled in with
confidence.
\item A dependency's engineering decisions are explained rather than worked
around: the crate that routes this part through its twin's peripheral module is
doing the right thing, and the README says why.
\item The conclusion is scoped to the ratio that was measured, and the operation
where the opposite conclusion holds is published in the same table.
\end{itemize}

\subsection*{Stretch goals}

\begin{itemize}
\item Add the services back to the register implementation one at a time,
namely the state check, the lock, the deadline and the error mapping, and plot
cost against services. The point at which the two lines meet is the real answer
to this chapter's question and nobody has published that curve.
\item Repeat the toggle measurement with the instruction cache disabled and with
the code placed in tightly coupled memory, which separates the abstraction's
cost from the memory system's. That is Chapter 19's subject approached from this
side.
\item Build the same four implementations at three optimisation levels and
publish the nine-cell table. Most of the disagreement in this field is people
comparing different cells of it.
\item Write the async framework example for this part properly, with a
peripheral that is not a pin, and offer it upstream. It is a small contribution
and it closes a real gap.
\item Measure charge per operation rather than cycles per operation with the
power instrument, which is the quantity that actually matters in Part II and
which nobody reports for abstraction layers at all.
\end{itemize}

\subsection*{Roadmap and next steps}

Chapter 18 takes the same instruments to a different question, namely what a
transform costs in two numeric formats with a numerical acceptance test, and
Chapter 19 explains the cache effect that this chapter's first run and median
are separated by. Chapter 20 is where the lock that the vendor layer takes stops
being decorative.

For going deeper, the community embedded engineering roadmap separates the
firmware path from the embedded Linux path and is explicit about weighting, and
the accompanying essay on bridging the gap makes the same argument in a page.
For the Rust direction this chapter opens, the standard progression is the
embedded working group's own book followed by the async framework's examples,
and the one caution worth repeating is the one this chapter is built on: check
that an example names your part and not its family. Appendix H lists the
courses; the free one whose code is under the most permissive licence in the
pool is the safest to build an exercise on, although its examples are not on
this board.

\subsection*{Portfolio evidence}

\begin{itemize}
\item A public repository containing four implementations of one operation that
provably return the same bytes, in two languages, with one build command each.
\item A results table whose every row carries its clock, optimisation level,
cache state, repetition count and statistic, which is the artefact this subject
currently lacks entirely.
\item A short written conclusion stating the ratio of overhead to work for two
operations at opposite ends of it, and saying plainly which of the two the
popular opinion is derived from.
\item A worked example for this exact part in a framework that had one only for
its twin, with the peripheral module routing explained.
\end{itemize}

\subsection*{Sources}

Normative references:
\begin{itemize}
\item Reference manual RM0455, for the I2C peripheral's transfer-configuration
register, the GPIO set and reset register, the cycle counter and the timers on
this part. Not RM0433, which documents its better-known sibling.
\item The STM32H7A3xI datasheet, for the clock limits that the cycle counts are
divided by.
\item The board user manual for MB1363, for the Arduino header mapping and the
free Zio pins used for the toggled pin and the marker.
\item The Cortex-M7 technical reference manual, for the cycle counter and the
cache behaviour that separates a first run from a median.
\end{itemize}

Reusable implementations:
\begin{itemize}
\item The only reproducible bench comparison of the vendor toggle call against a
direct register write, taken on a different member of this family at 550 MHz.\newline
\url{https://metebalci.com/blog/stm32h7-gpio-toggling/}
\item The Rust hardware crate for this family, 0BSD, whose feature for this part
routes to its twin's peripheral module and selects the right manual.\newline
\url{https://github.com/stm32-rs/stm32h7xx-hal}
\item The generated peripheral crates for this silicon, which carry a module for
this part.\newline
\url{https://github.com/stm32-rs/stm32-rs}
\item The async framework, Apache-2.0 or MIT, which names this exact package and
pin count and had no worked example for this part at the time of writing.\newline
\url{https://github.com/embassy-rs/embassy}
\item The real-time framework, which works here because it schedules on the
interrupt controller rather than on anything part-specific.\newline
\url{https://github.com/rtic-rs/rtic}
\item The flashing and debugging tool whose target manifest carries this exact
die.\newline
\url{https://github.com/probe-rs/probe-rs}
\item Oblaukhov's CMake package layer for this silicon family, MIT, which is the
practical way to build the vendor-layer implementation without its IDE.\newline
\url{https://github.com/ObKo/stm32-cmake}
\end{itemize}
